\section{Finite-record covariance energy: centering and collective memory}
This calculation concerns one fixed connectivity and one specified stationary
measure. Statistical variation of temporal windows is not independent sampling
of connectivity realizations. No decay model or infinite-record extrapolation
is fitted to the neuronal data below.

\subsection{Exact centering bias of one covariance block}
Write $z(t)=r(t)-m$ and
$C(\tau)=\mathbb E[z(t+\tau)z(t)^\top]$, with $C=C(0)$. For a block $B$
containing $L$ equally spaced samples, define
\[
 \widehat C_B=\frac1L\sum_{t\in B}z(t)z(t)^\top-
               \bar z_B\bar z_B^\top,
 \qquad \bar z_B=\frac1L\sum_{t\in B}z(t).
\]
This equals the covariance computed after subtracting the block's empirical
mean from the raw rates; knowledge of the true $m$ is not required. Stationarity
gives the exact expectation
\[
 M_B=\mathbb E[\bar z_B\bar z_B^\top]
     =\frac1{L^2}\sum_{t,u\in B}C(t-u),\qquad
 \boxed{\mathbb E\widehat C_B=C-M_B.}
\]
For a continuous block of duration $B$,
\[
 M_B=\frac1{B^2}\int_{-B}^{B}(B-|\tau|)C(\tau)\,\mathrm d\tau.
\]
The symmetric integral is a positive semidefinite matrix, even if the lagged
matrix is not symmetric. Moreover $0\preceq M_B\preceq C$ because
$C-M_B$ is an expected sample covariance.

The full-record variance numerator has the corresponding bias
\[
 \mathbb E\widehat\nu_T=\nu-\frac1N\operatorname{tr}M_T,
 \qquad \nu=N^{-1}\operatorname{tr}C.
\]
It is essential to distinguish the full record $T$, used in this numerator,
from the shorter blocks $B$, used in the energy denominator.

\subsection{Exact mean cross energy and its four-time term}
For blocks $B,D$ of equal duration, the separated-block estimator is
$\widehat{\mathcal E}_{BD}=N^{-1}\operatorname{tr}
(\widehat C_B\widehat C_D)$. Add and subtract the product of expectations:
\begin{align*}
 \mathbb E\widehat{\mathcal E}_{BD}
 &=\frac1N\operatorname{tr}[(C-M_B)(C-M_D)]+\chi_{BD},\\
 \chi_{BD}&=\frac1N\sum_{ij}
       \operatorname{Cov}(\widehat C_{B,ij},\widehat C_{D,ij}).
\end{align*}
At a large enough gap $\chi_{BD}$ may become negligible under mixing. The
centering bias remains. For independent equal-duration blocks,
\[
 \mathbb E\widehat{\mathcal E}_{BD}
 =\frac{\|C-M_B\|_F^2}{N}
 =\mathcal E-\frac{2\operatorname{tr}(CM_B)-\operatorname{tr}M_B^2}{N},
 \qquad\mathcal E=N^{-1}\operatorname{tr}C^2.
\]
Consequently increasing the gap alone does not produce an unbiased estimate.
Since $0\preceq M_B\preceq C$, the independent-block energy is biased downward.
If $\beta_B=\|M_B\|_F/\|C\|_F$, Cauchy--Schwarz gives
\[
 0\le\frac{\mathcal E-
       N^{-1}\|C-M_B\|_F^2}{\mathcal E}\le2\beta_B.
\]
Thus $\beta_B\le\eta/2$ is a sufficient centering-error criterion at relative
energy tolerance $\eta$, separate from the gap criterion for $\chi_{BD}$.

To show which correlation function controls $\chi$, introduce the centering
kernel $A_B(t,u)=\delta_{tu}/L-1/L^2$ on the block, zero outside. Then
\[
 \widehat C_{B,ij}=\sum_{t,u}A_B(t,u)z_i(t)z_j(u).
\]
Writing the fourth moment as its three pairings plus a fourth cumulant gives
\begin{align*}
 \chi_{BD}=\frac1N\sum_{t,u,s,v}A_B(t,u)A_D(s,v)
 \Big\{&\operatorname{tr}C(t-s)\operatorname{tr}C(u-v)\\
 &+\operatorname{tr}[C(t-v)C(u-s)]\\
 &+\sum_{ij}\operatorname{cum}[z_i(t),z_j(u),z_i(s),z_j(v)]\Big\}.
\end{align*}
This is exact on the discrete grid. It does not assume Gaussian rate
trajectories. Even a conditionally Gaussian input has nonlinear, non-Gaussian
rate outputs, so dropping the last term requires an additional approximation.
The formula makes explicit why a local two-time autocorrelation alone cannot
bound energy convergence: the estimator involves collective matrix pairings
and a four-time cumulant.

For bounded rates, a generic strong-mixing bound can make $|\chi_{BD}|$ at
most a constant times $N\alpha(G)$, where $G$ is the empty inter-block gap and
$\alpha$ is the mixing coefficient of the entire network. This is too crude
to set a useful record length without an $N$-uniform mixing estimate. We do
not replace that unknown coefficient by the decay time of an average neuron.

\subsection{An explicit collective-time criterion in a solvable illustration}
Before specializing the lag profile, suppose the matrix integrals and their
first absolute moments converge. With $S_C(0)=\int_{-\infty}^{\infty}C(\tau)
\,\mathrm d\tau$, the triangular window gives
\[
 M_B=B^{-1}S_C(0)+O(B^{-2}).
\]
Thus the general leading centering coefficients are
\begin{align*}
 \frac{\mathbb E\widehat{\mathcal E}_{BD}}{\mathcal E}
 &=1-\frac{4\tau_{\rm cent,E}}B+O(B^{-2}),\qquad
 \tau_{\rm cent,E}=\frac{\operatorname{tr}[C S_C(0)]}
                         {2\operatorname{tr}C^2},\\
 \frac{\mathbb E\widehat\nu_T}{\nu}
 &=1-\frac{2\tau_{\rm cent,\nu}}T+O(T^{-2}),\qquad
 \tau_{\rm cent,\nu}=\frac{\operatorname{tr}S_C(0)}{2\operatorname{tr}C}.
\end{align*}
These concern independent-block expectation bias. They are not automatically
the correlation time of a fourth-order pair process, nor the autocorrelation
time governing statistical fluctuations of a covariance-energy estimator.

The leading conditional theory can evaluate $\tau_{\rm cent,E}$ directly.
If $F(\omega_1,\omega_2)$ is its two-frequency energy density, then
\[
 \mathcal E=\iint F\,\frac{\mathrm d\omega_1\mathrm d\omega_2}{(2\pi)^2},
 \qquad J=\int F(\omega,0)\frac{\mathrm d\omega}{2\pi},\qquad
 \tau_{\rm cent,E}=\frac{J}{2\mathcal E}.
\]
This is a frequency \emph{slice}, not the width of the narrow anti-diagonal
resonance in the two-frequency plane. Likewise, before expansion the exact
independent-block centering prediction is obtained by multiplying the density
by $[1-a_B(\omega_1)][1-a_B(\omega_2)]$, where
$a_B(\omega)=\operatorname{sinc}^2(\omega B/2)$ uses $\operatorname{sinc}x=
\sin x/x$. Resolving $a_B$ needs a sufficiently fine frequency quadrature;
we do not multiply an unresolved narrow sinc on the coarse FFT grid.

At the archived coupling point, the leading conditional theory gives
$\tau_{\rm cent,E}=6.946$ (60:40), $6.889$ (80:20), while
$\tau_{\rm cent,\nu}=4.232,4.252$. At $B=500,T=4000$, the corresponding
first large-window estimates are a $5.56\%/5.51\%$ downward energy bias
and $5.13\%/5.09\%$ upward PR bias when ratio noise is ignored. These are
parameter-only estimates within the leading closure, not individual-network
corrections. A one-percent energy-centering target suggests $B\gtrsim2780$
in that approximation; gap and variance errors still require separate checks.

Consider commuting exponential covariance modes
\[
 C(\tau)=\sum_\alpha\lambda_\alpha
    e^{-|\tau|/\tau_\alpha}u_\alpha u_\alpha^\top.
\]
This is a stationary Gaussian illustration of estimator bias, not an assumed
EI covariance model. Performing the window integral gives
\begin{align*}
 M_B&=\sum_\alpha\lambda_\alpha b_\alpha(B)
               u_\alpha u_\alpha^\top,\\
 b_\alpha(B)&=\frac{2\tau_\alpha}{B}
       -\frac{2\tau_\alpha^2}{B^2}
          (1-e^{-B/\tau_\alpha}).
\end{align*}
For independent blocks,
\[
 \frac{\mathbb E\widehat{\mathcal E}_{BD}}{\mathcal E}
 =\frac{\sum_\alpha\lambda_\alpha^2[1-b_\alpha(B)]^2}
        {\sum_\alpha\lambda_\alpha^2}
 =1-\frac{4\tau_{\mathcal E}}B+O(B^{-2}),\quad
 \tau_{\mathcal E}=\frac{\sum_\alpha\lambda_\alpha^2\tau_\alpha}
                        {\sum_\alpha\lambda_\alpha^2}.
\]
The numerator instead has
\[
 \frac{\mathbb E\widehat\nu_T}{\nu}
 =1-\frac{2\tau_\nu}{T}+O(T^{-2}),\quad
 \tau_\nu=\frac{\sum_\alpha\lambda_\alpha\tau_\alpha}
                   {\sum_\alpha\lambda_\alpha}.
\]
Ignoring ratio fluctuations for this illustration, their combination gives
$\widehat d/d\simeq1+4\tau_{\mathcal E}/B-4\tau_\nu/T$. Shorter denominator
blocks can therefore bias PR upward even with an arbitrarily large block gap.
A sufficient asymptotic energy-centering rule in this illustration is
$B\gg4\tau_{\mathcal E}/\eta$, not merely $B\gg$ a local rate correlation time.

\subsection{Ratio fluctuations and practical diagnostics}
Let $b_\nu=\mathbb E\widehat\nu-\nu$ and
$b_E=\mathbb E\widehat E-\mathcal E$. A second-order delta expansion of
$\widehat d=\widehat\nu^2/\widehat E$ gives
\begin{align*}
 \frac{\mathbb E\widehat d-d}{d}
 \simeq{}&\frac{2b_\nu}{\nu}-\frac{b_E}{\mathcal E}
  +\frac{\operatorname{Var}\widehat\nu}{\nu^2}
  -\frac{2\operatorname{Cov}(\widehat\nu,\widehat E)}{\nu\mathcal E}
  +\frac{\operatorname{Var}\widehat E}{\mathcal E^2},\\
 \frac{\operatorname{Var}\widehat d}{d^2}
 \simeq{}&\frac{4\operatorname{Var}\widehat\nu}{\nu^2}
  -\frac{4\operatorname{Cov}(\widehat\nu,\widehat E)}{\nu\mathcal E}
  +\frac{\operatorname{Var}\widehat E}{\mathcal E^2}.
\end{align*}
Terms quadratic in the mean biases may be added when they are not small.
The cross-energy estimator is a product of covariance estimates, not the
squared Frobenius norm of the covariance used by its numerator. Its finite-T
ratio is not subject to the exact physical bound $d\le1$. An enormous value
signals an ill-conditioned ratio estimate, but alone does not prove a fixed
point, multistability or nonstationarity.

Our numerical tests vary block duration, empty gap and record duration
separately on one fixed network. They retain both fixed-number-of-block and
fixed-block-duration comparisons. The four-time quantity
\[
 K_{\rm outer}(\tau)=N^{-1}\langle[z(t)\cdot z(t+\tau)]^2\rangle-
                   N^{-1}\operatorname{tr}C^2
\]
is also measured as a diagnostic of instantaneous outer-product memory. It
does not by itself equal the full centered-block $\chi_{BD}$ above.
No empirical power law for required $T(N,\delta)$ is inferred from one network
and one coupling point. A critical scaling claim still needs matched duration
checks at several distances from onset and sizes, or a uniform theoretical
bound on the collective four-time integrals.

\subsection{A four-window mean-free repair}
The centering bias can be removed without knowing the true mean by using
independent copies. For mutually independent stationary draws $r^1,r^2$,
\[
 C=\frac12\mathbb E[(r^1-r^2)(r^1-r^2)^\top].
\]
For four independent draws this gives
\[
 \boxed{\mathcal E=\frac1{4N}\mathbb E\left[
       (r^1-r^2)^\top(r^3-r^4)\right]^2.}
\]
This is a mean-free independent-copy covariance identity, not a claim of a
new general statistical estimator. It suggests using four separated temporal
windows from the same stationary network. For two windows set
\begin{align*}
 \widehat C^{\rm diff}_{AB}
 &=\frac12\left(\widehat C_A+\widehat C_B\right)
    +\frac12(\bar r_A-\bar r_B)(\bar r_A-\bar r_B)^\top\\
 &=\frac1{2|A||B|}\sum_{t\in A,u\in B}
            [r(t)-r(u)][r(t)-r(u)]^\top.
\end{align*}
If the two windows are independent stationary records,
$\mathbb E\widehat C^{\rm diff}_{AB}=C$ exactly, regardless of their
internal temporal autocorrelation. If they are dependent, the expectation is
instead
\[
 \mathbb E\widehat C^{\rm diff}_{AB}
 =C-\frac12\left[
       \operatorname{Cov}(\bar r_A,\bar r_B)
      +\operatorname{Cov}(\bar r_B,\bar r_A)\right].
\]
Thus the mean-bias term now depends on the inter-window gap rather than an
unavoidable own-window $M_B$. For four independent stationary windows,
\[
 \widehat{\mathcal E}_{AB;CD}^{\rm diff}
 =N^{-1}\operatorname{tr}(\widehat C^{\rm diff}_{AB}
                         \widehat C^{\rm diff}_{CD})
\]
is exactly unbiased for $\mathcal E$. At finite gaps the remaining dependence
must still be assessed. Its variance can be larger for short windows, and
the ratio of averaged unbiased moments is not exactly unbiased for $d$.
Our implementation averages all three pairings of separated four-window
sets and does not count overlapping sets as independent observations.
Delete-one-window results are reported as sensitivity diagnostics; interpreting
them as temporal standard errors needs the mixing/independence approximation.

For $d_{AB}=\bar r_A-\bar r_B$, the compressed energy formula is
\begin{align*}
 4N\widehat{\mathcal E}_{AB;CD}^{\rm diff}
 ={}&\operatorname{tr}[(\widehat C_A+\widehat C_B)
                       (\widehat C_C+\widehat C_D)]\\
 &+d_{CD}^\top(\widehat C_A+\widehat C_B)d_{CD}
  +d_{AB}^\top(\widehat C_C+\widehat C_D)d_{AB}
  +(d_{AB}^\top d_{CD})^2.
\end{align*}
It allows the calculation from block means, covariance cross products and
quadratic forms without rebuilding a dense covariance for every quadruple.
The code verifies it against a direct dense product.
